\scanpage{4}

\uline{Lemme} Soit $C$ une courbe alg. non sing. complète \deleted{\unsure{alg|et}} irréd. de genre $g\geq 2$ sur $k$ alg. clos, $\sigma$ un automorphisme de $C$ induisant l'identité sur $J^0_{C/S}$. Alors $\sigma$ est l'identité.

En effet, le morphisme sur $J^1_{\deleted{C/S}}$ induit par $s$ est de la forme $T_a$, où $a$ est un pt de $J^0$ \unsure{rat|rat.} $/k$, égal à $\sigma x-x$ pour tt $x\in C(k)$. ($C$ étant considéré comme une \unsure{ptie|partie|plongé} de $J^1$).

Donc $\sigma(x)=x+a$, donc ce lemme équivaut au suivant (ds le cas particulier $S=\mathrm{Spec}(k)$)

\uline{Corollaire 2} \uline{Sous les conditions du th}, \deleted{Toute} \inserted{Une} section $s$ de $J^0_{C/S}$ sur $S$ telle que la translation $T_s$ de $J^1_{C/S}$ invarie $C$ (plongée dans $J^1_{C/S}$) est \deleted{l'identité} \inserted{section nulle}

En effet, il résulte de la formule des Traces que le ``\unsure{nb|nb.} des pts. fixes'' \deleted{de} $\sigma$ est égal à
\[
\deleted{1-2g+1=2-2g}\quad \inserted{-1+2g-1=2(g-1)>0}
\]
(puisque $\sigma$ opère trivialement dans la jacobienne $J^0$, donc sa trace y est $2g$] et si $x$ est un pt fixe, on a $a=\sigma x-x=0$

% archive sheet number: 3
